% ============================================================================
%  Actividad 4 · Pre-validacion con evaluadores prototipo e iteracion
%
%  Propiedad de US-UX-07. Responde al criterio 6 de la historia: pre-validacion
%  sobre capturas reales del prototipo, con al menos una iteracion documentada
%  como hallazgo, cambio y version, con el antes y el despues lado a lado.
%
%  La ronda produjo tres hallazgos y las tres iteraciones estan documentadas.
%  Las cifras de contraste de la segunda salen de correr design/contraste.py:
%  no se estiman ni se copian de la planeacion.
%
%  Las dos figuras del par salen de figuras/a4/antes/ y figuras/a4/despues/, y
%  se emparejan por igualdad de nombre de archivo.
% ============================================================================

% ============================================================================
\uxsection{Pre-validación e iteración}

Antes de poner el prototipo frente a personas conviene gastar una ronda barata que localice los
defectos que se ven a simple vista. Esta sección documenta esa ronda: quién la ejecutó, sobre qué
material, con qué tareas y qué cambió el prototipo como consecuencia.

\textbf{Los ocho evaluadores son prototipos sintéticos.} Cada uno está condicionado por una de las
ocho personas construidas en la Actividad 1, siguiendo el mismo protocolo de condicionamiento que
sostuvo el card sorting de la Actividad 3. No son participantes humanos, sus respuestas no son datos
de campo y no se mezclan con ninguna cifra proveniente de instrumentos aplicados a personas reales.
Son predicciones plausibles de cómo reaccionaría cada perfil, y su utilidad está en detectar
ambigüedad y defectos de presentación antes de gastar la única ronda con personas que el proyecto
tiene, que es la prueba de usabilidad de la Actividad 5.

\subsection{Material evaluado}

El material fueron las \textbf{siete capturas reales del prototipo}, tomadas con el protocolo de la
sección de método: resolución fija, modo claro, idioma español y una sesión por pantalla con el rol
correspondiente. No se evaluaron maquetas ni pantallas dibujadas para la ocasión, que es la
condición que la actividad pone a esta ronda: un hallazgo sobre un dibujo no dice nada sobre el
producto.

\subsection{Los ocho evaluadores}

Se conservan los mismos perfiles del ejercicio de la Actividad 3, con la misma numeración, para que
los hallazgos de las dos rondas puedan compararse entre sí.

\begin{uxtabla}{|L{2.4cm}|L{3.4cm}|Y|}{Los ocho evaluadores prototipo y lo que cada uno vigila}
  \uxheadrow \thd{Evaluador} & \thd{Perfil} & \thd{Qué mira primero en una pantalla} \\
  P1 Laura & Operativo & Si puede localizar una cifra y comprobarla sin preguntar a nadie \\
  P2 Diego & Analista & Si entiende el dato antes de usarlo y si puede repetir la consulta mañana \\
  P3 Roberto & Propietario de datos & Si las definiciones oficiales y la autorización de acceso a sus bases están a la vista \\
  P4 Elena & Auditoría & Si hay rastro y evidencia: de dónde sale la cifra y quién la consultó \\
  P5 Jorge & Ingeniería de datos & Si el estado de las cargas y la estructura del dato son visibles \\
  P6 Arturo & Directivo & Si lo que necesita antes de una junta está en la primera pantalla y con respaldo \\
  P7 Mariana & Administración & Si cuentas, permisos y continuidad del servicio se administran desde un solo lugar \\
  P8 Ximena & Integración & Si el contrato del dato y la vía de salida están documentados \\
\end{uxtabla}

\subsection{Las dos tareas heredadas de la Actividad 3}

La prueba de árbol de la actividad anterior dejó dos preguntas abiertas, y esta ronda las retoma
sobre la arquitectura ya revisada. Las dos se formulan igual que entonces, para que la comparación
tenga sentido.

\begin{uxlist}
  \lead{Tarea 1, el indicador de riesgo}{En la Actividad 3, consultar un indicador de riesgo obtuvo
    un solo acierto de ocho contra la ruta esperada en gobierno del dato, y seis de los ocho
    evaluadores lo buscaron en exploración. La arquitectura se corrigió y los tableros pasaron a
    exploración. Sobre las capturas revisadas, la ruta que los evaluadores señalan coincide con la
    que el prototipo implementa: la subrama de tableros e indicadores aparece en el segundo nivel de
    exploración, visible en la barra lateral de la captura de esa pantalla.}
  \lead{Tarea 2, la bitácora de accesos}{Revisar quién autorizó un acceso hizo titubear a los ocho
    evaluadores, la única tarea con duda unánime, y dos de ellos la buscaron directamente en
    gobierno. La arquitectura mantuvo la bitácora en administración con un acceso cruzado desde
    gobierno del dato. Sobre las capturas, el acceso cruzado está presente y la duda se reduce, pero
    no desaparece: sigue habiendo un momento de decisión entre las dos categorías.}
\end{uxlist}

Las dos observaciones se registran como \textbf{orientación y no como medición}. Con evaluadores
sintéticos no se reclama un porcentaje de acierto comparable con el de la prueba de árbol original;
lo que aportan es la dirección del hallazgo. La medición que cierra estas dos preguntas es la de la
Actividad 5, con participantes reales.

% ============================================================================
\section{Primera iteración: la franja de alcance}

De la ronda salieron \textbf{tres hallazgos} y los tres produjeron un cambio. El primero se aplicó
de inmediato porque cumplía las tres condiciones que hacen útil una iteración en esta etapa:
aparecía en todas las pantallas, era barato y su efecto se comprueba mirando la imagen. Los otros
dos quedaron registrados con su motivo de aplazamiento y se cerraron después, dentro de la misma
semana; las secciones siguientes los documentan con el mismo esquema de hallazgo, cambio y versión.

\subsection{Hallazgo}

\textbf{La franja que declara el alcance del prototipo se leía como una tarjeta decorativa y no como
una condición del sistema.} El aviso ---\enquote{Prototipo de alta fidelidad de Karisma Data con
datos sintéticos. No está conectado a sistemas reales de ninguna institución.}--- aparecía dentro de
una caja estrecha pegada a la esquina superior izquierda del contenido, con su texto partido en dos
líneas y el resto de la banda vacío.

El hallazgo lo levantan con más fuerza los dos perfiles que viven de la procedencia del dato. Elena,
desde auditoría, mira primero de dónde sale una cifra y quién la consultó; Roberto responde por las
bases que administra. Para ambos, un aviso que se lee como adorno es un aviso que no cumple su
función. Arturo, que revisa cifras antes de una junta, aporta la consecuencia práctica: si el aviso
parece decorativo, una captura de esa pantalla puede circular en una presentación sin que nadie
registre que los datos son sintéticos.

La causa era técnica y se midió en el navegador en lugar de suponerse. La franja se compone como un
párrafo, y el sistema de diseño aplica a todo párrafo una medida de lectura de sesenta y ocho
caracteres para que la prosa no se extienda más allá de lo cómodo. La franja heredaba esa medida:
resultaba de \textbf{455 píxeles} de ancho dentro de una columna de \textbf{1193}. La regla es
correcta para prosa y no le corresponde a una declaración de sistema, que no se lee de corrido sino
de un vistazo.

\subsection{Cambio}

La franja deja de heredar la medida de lectura y ocupa el ancho de su columna. El cambio se aplica
en las tres composiciones que la montan y no en el componente, porque el componente pertenece a otra
historia de usuario en esta misma semana y modificarlo habría cruzado un límite de propiedad. La
excepción se declara en el mismo lugar donde se aplica, con la medición que la motiva.

El texto se mantiene en una sola línea y la banda gana altura útil: la caja pasa de \textbf{48 a 33
píxeles} de alto, de modo que la declaración ocupa \emph{menos} espacio vertical y \emph{más} ancho.
Es el resultado buscado: un aviso que se percibe como una franja del sistema y no como un elemento
más del contenido.

\subsection{Versión}

\begin{uxtabla}{|L{3.0cm}|Y|}{Registro de la iteración}
  \uxheadrow \thd{Elemento} & \thd{Contenido} \\
  Hallazgo & La franja de alcance se lee como tarjeta decorativa y no como condición del sistema \\
  Origen & Ronda de pre-validación sobre las siete capturas; levantado por los perfiles de auditoría, propiedad del dato y dirección \\
  Medición & Ancho de 455 píxeles dentro de una columna de 1193; alto de 48 píxeles en dos líneas \\
  Cambio & La franja deja de heredar la medida de lectura y ocupa el ancho de la columna \\
  Verificación & Ancho de 1193 píxeles, igual al de la columna; alto de 33 píxeles en una línea \\
  Alcance & Las siete pantallas, porque la franja es común a todas \\
\end{uxtabla}

\subsection{El antes y el después}

Las dos capturas siguientes corresponden a la misma pantalla, con la misma sesión, la misma
resolución y el mismo protocolo. La única diferencia entre ellas es el cambio descrito.

\figuraux{figuras/a4/antes/2_exploracion_normal.png}{Antes. La franja de alcance ocupa 455 píxeles de
una columna de 1193 y su texto se parte en dos líneas, de modo que el aviso se percibe como una
tarjeta suelta en la esquina superior izquierda.}{fig:a4-antes}

\figuraux{figuras/a4/despues/2_exploracion_normal.png}{Después. La franja ocupa el ancho de la
columna y su texto cabe en una línea. La declaración se lee como una condición de toda la pantalla,
que es su función.}{fig:a4-despues}

La captura del antes se tomó \textbf{antes} de aplicar ningún hallazgo, y esa precedencia no es un
detalle de procedimiento: una vez modificada la interfaz, el estado anterior no se puede reconstruir
de memoria y la iteración deja de ser verificable. Por eso se capturaron las siete pantallas en la
primera pasada y no solo la que se preveía cambiar: al momento de capturar todavía no se sabía cuál
cambiaría.

% ============================================================================
\section{Segunda iteración: la paleta bajo el listón de contraste}

El segundo hallazgo de la ronda quedó registrado el mismo día y se cerró después, cuando dejó de
depender de una decisión ajena. Es la iteración que más consecuencias tuvo, y ninguna de ellas se
ve en una captura: se ven en números.

\subsection{Hallazgo}

\textbf{La primera impresión del prototipo ocurría en una paleta que el propio documento declaraba
no haber verificado.} En un perfil de navegador nuevo el modo de color sigue la preferencia del
sistema operativo, de modo que quien tuviera su equipo en oscuro veía el portal en oscuro; y la guía
de estilos declaraba el modo oscuro fuera de alcance justamente por no tener calculada su matriz de
contraste. Un aviso de este tipo no se resuelve escribiendo una advertencia: se resuelve calculando
la matriz que falta o retirando la paleta.

Al mismo hallazgo se sumó una exigencia de identidad: el portal tenía que poder mostrarse con la
paleta institucional del archivo de diseño sin renunciar a ninguna de las garantías que la guía ya
publicaba. Las dos cosas se resolvieron con el mismo trabajo, porque el listón es el mismo.

\subsection{Medición}

Antes de aceptar ningún color se corrió sobre la paleta candidata el mismo cálculo que sostiene la
matriz de la guía de estilos, en los dos modos, y con las simulaciones de las tres dicromacias. La
paleta no pasó, y los defectos aparecieron con nombre y con cifra:

\begin{uxlist}
  \lead{El color de atención no alcanzaba el listón de texto}{Daba \textbf{3{,}65:1} sobre el suelo
    claro, por debajo del 4{,}5:1 que la guía exige de todo token que informa.}
  \lead{El canal informativo se confundía con el de confirmación}{El color de acción y el de éxito
    quedaban a \textbf{6{,}7} de distancia perceptual bajo tritanopia, muy por debajo del piso de
    \textbf{13{,}6} que el tema de omisión sostiene. Para quien no distingue azul de amarillo, un
    aviso informativo y una confirmación se veían igual.}
  \lead{El color de error se acercaba demasiado al de éxito bajo protanopia}{Este tercer defecto no
    estaba previsto y apareció al medir: \textbf{10{,}1} en modo claro contra un piso de 13{,}6 y
    \textbf{13{,}1} en modo oscuro contra un piso de 21{,}5. La revisión anterior solo había mirado
    los pares del canal informativo, y este par vive en otro.}
\end{uxlist}

\subsection{Cambio}

Tres correcciones, cada una con su regla escrita antes de aplicarla. El color de atención se oscurece
un 12\,\% conservando matiz y saturación, que es el mínimo que cruza el listón: a un 10\,\% todavía
daba 4{,}40:1. El canal informativo sale del octeto de la marca, porque dentro de él no hay ningún
matiz que se separe del de confirmación bajo tritanopia: en claro se resuelve con un azul profundo
y en oscuro con un violeta, ya que bajo esa dicromacia el azul y el verde colapsan. El color de
error se oscurece un 20\,\% en claro y un 12\,\% en oscuro, con la misma regla de conservación de
matiz.

\subsection{Versión}

\begin{uxtabla}{|L{3.0cm}|Y|}{Registro de la segunda iteración}
  \uxheadrow \thd{Elemento} & \thd{Contenido} \\
  Hallazgo & La primera impresión ocurría en una paleta declarada sin verificar, y la paleta institucional no había pasado por el cálculo \\
  Origen & Ronda de pre-validación sobre las siete capturas, más la exigencia de identidad del archivo de diseño \\
  Medición & Atención en 3{,}65:1 sobre el suelo claro; informativo y confirmación a 6{,}7 bajo tritanopia, con piso 13{,}6; error y confirmación a 10{,}1 en claro y 13{,}1 en oscuro bajo protanopia \\
  Cambio & Atención oscurecida 12\,\%; canal informativo fuera del octeto, azul profundo en claro y violeta en oscuro; error oscurecido 20\,\% en claro y 12\,\% en oscuro \\
  Verificación & Atención en 4{,}54:1; error en 7{,}46:1 y 5{,}13:1 sobre su suelo; peor par bajo dicromacia de 14{,}5 en claro y 22{,}6 en oscuro, por encima de los pisos de 13{,}6 y 21{,}5; \textbf{cero} tokens que informan por debajo de su listón en las cuatro combinaciones de tema y modo; las seis series de datos por encima de 3:1 en las cuatro, con un mínimo de 4{,}64:1 \\
  Alcance & Los diecisiete tokens de color del sistema, en dos temas y dos modos \\
  Versión & El cambio se publica dentro de la versión v2.0 del sistema de tokens del portal, fechada el 16 de agosto de 2026. La guía de estilos de este documento conserva su v1.0: es otra cadena y no se movió \\
\end{uxtabla}

Esta iteración no tiene par de capturas y no lo necesita, porque su evidencia es de otra naturaleza.
Un lector puede volver a correr el cálculo sobre la definición de los tokens y obtener las mismas
cifras; una fotografía de la pantalla, en cambio, no distingue 3{,}65:1 de 4{,}54:1. La paleta
completa, con el origen de cada valor, se publica en la sección dedicada al tema institucional y a
los flujos de tarea.

El hallazgo de partida queda cerrado con esto: el modo oscuro dejó de ser una paleta sin verificar.
El protocolo de captura sigue fijando el modo claro y el tema de omisión, y ahora por una razón
distinta ---mantener comparables las siete capturas y el par de la primera iteración--- y no por
falta de verificación.

% ============================================================================
\section{Tercera iteración: la pantalla que declaraba en lugar de mostrar}

\subsection{Hallazgo}

\textbf{Varios evaluadores señalaron la pantalla de exploración como la más pobre del conjunto, y
tenían razón.} Era la única de las siete que declaraba sus capacidades en lugar de ejercerlas:
enumeraba lo que alojaría y nombraba la historia de usuario que lo entrega. El estado era honesto y
seguía siendo el punto débil del recorrido, porque la categoría que el card sorting de la Actividad
3 identificó como el centro de la tarea de consulta era precisamente la que no se podía usar.

El aplazamiento inicial tenía un motivo y se escribió: construir el catálogo temático era trabajo de
producto, no de presentación, y la alternativa barata ---rellenar la pantalla con contenido de
ejemplo--- habría presentado como producto algo que no lo era.

\subsection{Cambio}

La pantalla dejó de montar el componente de andamiaje y pasó a componer el catálogo desde la misma
consulta que ya alimentaba el diccionario de campos de gobierno del dato. Con ella entraron sus
cuatro estados no felices, que la sección de prototipos detalla: vacío con explicación, carga con
esqueleto de la altura de las filas reales, error que conserva el término escrito, y sin permiso
sobre las dos continuaciones del módulo que sí exigen perfil.

\subsection{Versión}

\begin{uxtabla}{|L{3.0cm}|Y|}{Registro de la tercera iteración}
  \uxheadrow \thd{Elemento} & \thd{Contenido} \\
  Hallazgo & La pantalla de exploración declaraba sus capacidades en lugar de ejercerlas \\
  Origen & Ronda de pre-validación sobre las siete capturas; señalado por varios evaluadores como el punto más débil del recorrido \\
  Medición & Una de siete pantallas sin contenido propio; tres de sus cuatro estados no felices resueltos por componentes de otras historias y el cuarto solo especificado \\
  Cambio & Buscador, conteos por dominio y resultados compuestos desde la búsqueda del catálogo, con los cuatro estados construidos dentro de la pantalla \\
  Verificación & Cero apariciones del componente de andamiaje en la pantalla y once casos de prueba automática que montan sus estados; la matriz de estados no felices pasa de 6 a 10 celdas construidas en esta actividad \\
  Alcance & La pantalla de exploración y las dos continuaciones gobernadas de su módulo \\
  Versión & Publicada en el contrato de navegación como \emph{navegable con datos de ejemplo}, que es el estado que la tabla de alcance declara \\
\end{uxtabla}

\textbf{Esta iteración tampoco se recaptura, y la decisión es deliberada.} La imagen del estado
anterior existe y es la figura del prototipo 2; volver a fotografiar esa pantalla habría producido
una imagen mejor y habría destruido el par de la primera iteración, que compara exactamente la misma
pantalla antes y después del cambio de la franja. Esa pareja no se puede volver a tomar, porque el
estado del antes ya no existe. Entre una figura más atractiva y una evidencia irrepetible, se
conserva la evidencia.

% ============================================================================
\section{Lo que queda abierto}

Tres cosas, dichas aquí para que no haya que deducirlas.

\begin{uxlist}
  \lead{El catálogo construido no tiene captura en este documento}{Su estado entregado se describe
    en la sección de prototipos y se verifica con pruebas automáticas, no con una imagen. La razón
    es la del párrafo anterior: la única captura de esa ruta sostiene un par que no se puede
    rehacer.}
  \lead{Las dos preguntas de la prueba de árbol siguen sin medición humana}{La ruta del indicador
    de riesgo y la del acceso cruzado a la bitácora se revisaron con evaluadores prototipo, que dan
    dirección y no porcentaje. La medición corresponde a la prueba de usabilidad de la Actividad 5.}
  \lead{El cálculo de contraste cubre color y no cubre densidad}{La matriz verifica legibilidad por
    contraste en las cuatro combinaciones de tema y modo. La comodidad de lectura a la altura de
    fila del producto es una cuestión distinta y se mide con personas, no con un cálculo.}
\end{uxlist}
